\subsection{Boosting-Based Tree Ensembles}
\label{subsec:rw_boosting}

Boosting constructs an ensemble sequentially so that later learners focus on
errors made by earlier learners. This makes boosting well suited to structured
tabular prediction because nonlinear effects and feature interactions can be
captured without requiring a single highly complex tree. The three models used
in this study represent different generations of boosting-based tree methods:
AdaBoost provides a classical boosting baseline, while XGBoost and LightGBM are
modern gradient-boosted decision-tree systems designed for stronger predictive
performance and computational efficiency.

\subsection{AdaBoost}
\label{subsec:rw_adaboost}

AdaBoost is one of the foundational boosting algorithms. It updates observation
weights after each boosting round, increasing emphasis on samples that were
previously misclassified, and combines the weak learners into a weighted final
classifier \cite{freund1997adaboost}. In this study, AdaBoost serves as the
classical baseline and uses decision stumps as weak learners. Keeping the base
learners simple provides a clear comparison with the deeper and more heavily
regularized gradient-boosting systems used by XGBoost and LightGBM.

\subsection{XGBoost}
\label{subsec:rw_xgboost}

XGBoost extends gradient tree boosting with both statistical regularization and
systems-level optimizations for scalable learning. Its design includes a
regularized objective, sparsity-aware split handling, approximate tree learning,
and efficient parallel computation \cite{chen2016xgboost}. It also exposes
controls such as tree depth, shrinkage, row subsampling, column subsampling, and
regularization. These controls are useful for the present study because the
three healthcare datasets differ substantially in sample size, dimensionality,
and class prevalence, while the same fixed model configuration must remain
reproducible across datasets.

\subsection{LightGBM}
\label{subsec:rw_lightgbm}

LightGBM is a gradient-boosted decision-tree framework designed for high
computational efficiency. Ke et al. introduced Gradient-based One-Side Sampling
(GOSS) and Exclusive Feature Bundling (EFB) to reduce computational cost while
preserving predictive quality on large and high-dimensional datasets
\cite{ke2017lightgbm}. LightGBM also grows trees leaf-wise, allowing it to fit
complex nonlinear relationships efficiently. Because leaf-wise growth can
increase model complexity, parameters such as the number of leaves, maximum
depth, minimum child samples, learning rate, subsampling, and regularization are
controlled explicitly in the experimental configuration.

\subsection{Explainability and Stability of Tree Models}
\label{subsec:rw_shap_stability}

Predictive performance alone does not show which variables drive model outputs.
SHAP (SHapley Additive exPlanations) provides additive feature attributions for
individual predictions and establishes desirable consistency properties for
feature attribution \cite{lundberg2017shap}. Tree-specific SHAP algorithms make
these explanations computationally practical for tree ensembles, and prior work
has demonstrated their use in medical prediction settings
\cite{lundberg2020treeexplainer}. This motivates evaluating not only predictive
performance across cross-validation folds but also whether feature-importance
rankings remain consistent as the training sample changes.

However, healthcare prediction studies are difficult to compare directly when
they use different datasets, preprocessing choices, class-balance strategies,
train--test partitions, or evaluation metrics. A reported advantage for one
model may therefore reflect the experimental setup rather than a stable property
of the algorithm. Explainability studies also commonly report a single global
ranking, which does not establish whether that ranking persists when the training
sample changes. These limitations leave two related questions open: whether
model performance is stable across folds and datasets, and whether the features
used to explain that performance are stable under the same changes.

\subsection{Research Gap and Motivation}
\label{subsec:rw_gap}

The gap addressed by this study is therefore methodological rather than a claim
that one boosting algorithm is universally best. AdaBoost, XGBoost, and LightGBM
are evaluated under one documented preprocessing, fold-generation, metric, and
hold-out protocol across three healthcare datasets. The same persisted folds are
used for model evaluation and SHAP analysis, while preprocessing is fitted only
on each fold's training rows. The study consequently separates average
predictive strength (RQ1), fold-to-fold predictive variability (RQ2), and
fold-to-fold explanation variability (RQ3). This controlled design does not
eliminate differences between the datasets, but it makes those differences
explicit rather than conflating them with different evaluation procedures.
